% Part: first-order-logic
% Chapter: syntax-and-semantics
% Section: free-vars-sentences

\documentclass[../../../include/open-logic-section]{subfiles}

\begin{document}

\olfileid{fol}{syn}{fvs}

\olsection{मुक्त \printtoken{P}{variable} आणि \printtoken{P}{sentence}}

\begin{defn}[!!a{variable} च्या मुक्त आस्थिती]
\ollabel{defn:free-occ}
!!a{formula} मधील !!a{variable} च्या \emph{मुक्त} आस्थिती पुढीलप्रमाणे
विगमनाने परिभाषित केल्या जातात:
\begin{enumerate}
\item \indcase*{!A}{$!A$ आण्विक आहे}{$\indfrm$ मधील सर्व
  !!{variable}-आस्थिती मुक्त असतात.}

\tagitem{prvNot}{\indcase{!A}{\lnot !B}{$\indfrm$ मधील मुक्त
  !!{variable}-आस्थिती म्हणजे नेमक्या $!B$ मधील मुक्त आस्थिती.}}{}

\item \indcase{!A}{(!B \ast !C)}{$\indfrm$ मधील मुक्त
  !!{variable}-आस्थिती म्हणजे $!B$ मधील आस्थिती आणि त्यांच्यासोबत
  $!C$ मधील आस्थिती.}

\tagitem{prvAll}{\indcase{!A}{\lforall[x][!B]}{$\indfrm$ मधील मुक्त
  !!{variable}-आस्थिती म्हणजे $!B$ मधील $x$ च्या आस्थिती सोडून इतर
  सर्व मुक्त आस्थिती.}}{}

\tagitem{prvEx}{\indcase{!A}{\lexists[x][!B]}{$\indfrm$ मधील मुक्त
  !!{variable}-आस्थिती म्हणजे $!B$ मधील $x$ च्या आस्थिती सोडून इतर
  सर्व मुक्त आस्थिती.}}{}
\end{enumerate}
\end{defn}

\begin{defn}[बद्ध चरे]
सूत्र~$!A$ मधील !!a{variable} ची आस्थिती मुक्त नसेल, तर ती
\emph{बद्ध} असते.
\end{defn}

\begin{prob}
\olref[fol][syn][fvs]{defn:free-occ} च्या धर्तीवर बद्ध चर-आस्थितींची
विगमनाधारित व्याख्या द्या.
\end{prob}

\begin{defn}[व्याप्ती]
\iftag{prvAll}{सूत्र~$!A$ मध्ये $\lforall[x][!B]$ ची एखाद्या उपसूत्राची
  आस्थिती असेल, तर $!A$ मधील $!B$ च्या संबंधित आस्थितीला
  $\lforall[x]$ च्या संबंधित आस्थितीची \emph{व्याप्ती} म्हणतात.
  \iftag{prvEx}{$\lexists[x]$ साठीही त्याचप्रमाणे.}{}}{सूत्र~$!A$
  मध्ये $\lexists[x][!B]$ ची एखाद्या उपसूत्राची आस्थिती असेल, तर
  $!A$ मधील $!B$ च्या संबंधित आस्थितीला $\lexists[x]$ च्या संबंधित
  आस्थितीची \emph{व्याप्ती} म्हणतात.}

$!A$ मधील संख्यापकाच्या
\iftag{prvAll}{$\lforall[x]$\iftag{prvEx}{ किंवा
    $\lexists[x]$}{}}{$\lexists[x]$} या आस्थितीची व्याप्ती $!B$ असेल,
तर $!B$ मधील $x$ च्या मुक्त आस्थिती
\iftag{prvAll}{$\lforall[x][!B]$\iftag{prvEx}{ आणि
$\lexists[x][!B]$}{}}{$\lexists[x][!B]$} मध्ये बद्ध असतात. या आस्थिती
त्या उल्लेख केलेल्या संख्यापक-आस्थितीने \emph{बद्ध केल्या आहेत} असे
आपण म्हणतो.
\end{defn}

\begin{ex}
पुढील सूत्राचा विचार करा:
\[
\lexists[\Obj v_0][\underbrace{\Atom{\Obj A^2_0}{\Obj v_0,\Obj v_1}}_{!B}]
\]
$!B$ हे $\lexists[\Obj v_0]$ ची व्याप्ती दर्शवते. हा संख्यापक $!B$
मधील $\Obj v_0$ ची आस्थिती बद्ध करतो; पण $\Obj v_1$ ची आस्थिती बद्ध
करत नाही. म्हणून या बाबतीत $\Obj v_1$ हे मुक्त चर आहे.


अधिक गुंतागुंतीच्या !!{formula}~$!A$ मध्ये हे कसे कार्य करते ते आता
पाहू शकतो:
\[
\lforall[\Obj v_0][\underbrace{(\Atom{\Obj A^1_0}{\Obj v_0} \lif
    \Atom{\Obj A^2_0}{\Obj v_0, \Obj v_1})}_{!B}] \lif \lexists[\Obj
  v_1][\underbrace{(\Atom{\Obj A^2_1}{\Obj v_0, \Obj v_1} \lor \lforall[\Obj v_0][\overbrace{\lnot \Atom{\Obj A^1_1}{\Obj v_0}}^{!D}])}_{!C}]
\]
$!B$ ही पहिल्या $\lforall[\Obj v_0]$ ची व्याप्ती, $!C$ ही
$\lexists[\Obj v_1]$ ची व्याप्ती आणि $!D$ ही दुसऱ्या
$\lforall[\Obj v_0]$ ची व्याप्ती आहे. पहिले $\lforall[\Obj v_0]$ हे
$!B$ मधील $\Obj v_0$ च्या आस्थिती बद्ध करते, $\lexists[\Obj v_1]$ हे
$!C$ मधील $\Obj v_1$ ची आस्थिती बद्ध करते आणि दुसरे
$\lforall[\Obj v_0]$ हे $!D$ मधील $\Obj v_0$ ची आस्थिती बद्ध करते.
$\Obj v_1$ ची पहिली आस्थिती आणि $\Obj v_0$ ची चौथी आस्थिती $!A$ मध्ये
मुक्त आहेत. $\Obj v_0$ ची शेवटची आस्थिती $!D$ मध्ये मुक्त, पण $!C$
आणि~$!A$ मध्ये बद्ध आहे.
\end{ex}

\begin{defn}[वाक्य]
!!^a{formula}~$!A$ मध्ये !!{variable}s च्या मुक्त आस्थिती नसतील तेव्हा
आणि केवळ तेव्हाच ते \article{sentence} \emph{!!{sentence}} असते.
\end{defn}

% add examples!

\end{document}
